\documentclass[12pt,a4paper]{article}
\usepackage[margin=2.5cm]{geometry}
\usepackage{times}
\usepackage{url}

\begin{document}

\begin{center}
{\Large\bfseries SciFig-Eval: A Multilingual Benchmark and\\
Behavioural Framework for Evaluating Vision-Language\\
Models on Scientific Figures}

\bigskip

{\large Paul Osemudiame Oamen}

\medskip

{\normalsize Department of Computing Science\\
University of Aberdeen\\
MSc Artificial Intelligence, 2026}

\bigskip

{\large\bfseries Abstract}
\end{center}

\medskip

\noindent
Vision-language models are rapidly becoming integral to the scientific
research workflow, supporting tasks from automated literature review
and paper evaluation to research assistance and data interpretation.
Central to these capabilities is the ability to comprehend scientific
figures, which typically encode the core quantitative arguments of
published research. As adoption accelerates, it becomes important to
investigate how well these models comprehend such visuals, and beyond
simple task accuracy, to distill their behavioural profiles
(encompassing epistemic honesty, adversarial robustness, and contextual
reasoning) that current benchmarks leave largely unexamined.

To address this gap, we introduce SciFig-Eval, a multilingual
evaluation corpus of 1,005 scientific figures drawn from over 300
research papers across four languages (English, Bulgarian, Chinese,
German) with 1,411 expert annotations, 630 capability questions, and
adversarial probes grounded in the misinformation effect, anchoring
bias, and sycophancy, scored using an MQM framework. We further propose
the Admittance-Resistance-Inductance (A-R-I) framework, which
decomposes figure-reading trustworthiness into three behavioural axes
scored through a dual-judge pipeline validated against human annotators
(ICC\,=\,.91).

Evaluating thirteen frontier models, spanning open-weight architectures
from 4B to 235B parameters alongside leading closed-source systems, we
find that behavioural profiles diverge sharply from accuracy rankings.
GPT-5.2 leads on description quality but scores only .56 on resistance
to false premises, the Gemma family fabricates with near-zero admittance
regardless of scale, numerical precision and visual counting account for
the majority of errors, and description and comprehension prove to be
partially dissociated. These findings demonstrate that trustworthy
deployment requires behavioural profiling in addition to accuracy
measurements.

\end{document}
